%================================================================
% Poglavje 1: Uvod
%================================================================
\chapter{Uvod}
\label{ch:uvod}

Problem optimizacije naložbenih portfeljev (v nadaljevanju ONP) predstavlja enega temeljnih izzivov v kvantitativnih financah. Markowitzev model~\cite{markowitz1952} išče optimalne uteži za nek nabor finančnih sredstev, s katerimi dosežemo najboljše ravnovesje med pričakovanim donosom in tveganjem. Kljub svoji eleganci ima v osnovni obliki dve izraziti pomanjkljivosti, ki resno omejujeta njegovo uporabnost v praksi in ju celovito naslavljamo v tem magistrskem delu.

Prva pomanjkljivost je zanašanje na pretekle podatke pri oceni prihodnjih pričakovanih donosov ($\vec{\mu}$) in kovariance ($\vec{\Sigma}$). Ker optimizacijski algoritmi napake v oceni $\vec{\mu}$ pogosto ojačajo, je zanesljiva napoved ključnega pomena. Ta problem v nalogi rešujemo z uporabo transformerskih modelov~\cite{vaswani2017}, natančneje arhitekture MASTER~\cite{li2024master}, ki iz iste deljene predstavitve hkrati napoveduje tako pričakovane donose kot matriko tveganja.

Druga pomanjkljivost je nerealistična predpostavka o razpršitvi premoženja po vseh razpoložljivih sredstvih. V praksi vlagatelji želijo zmanjšati transakcijske stroške in likvidnostne omejitve, kar zahteva portfelj z omejenim številom različnih sredstev. Vpeljava te kardinalne omejitve (največ $K$ izbranih sredstev) pretvori prvotni problem v NP-težkega~\cite{chang2000}. Za iskanje optimalnih uteži zato uporabljamo uveljavljene metahevristične postopke, kot so rojenje delcev, simulirano ohlajanje in genetski algoritmi.

ONP ni zgolj teoretičen izziv, temveč vsakodnevna naloga upraviteljev premoženja --- od pokojninskih in investicijskih skladov prek samodejnih svetovalcev (angl.\ \emph{robo-advisors}) do posameznih vlagateljev ---, ki upravljajo znatna sredstva in pri katerih že majhna izboljšava tvegano prilagojenega donosa pomeni veliko absolutno vrednost. Prav zato sta kakovost ocen momentov donosov in upoštevanje realističnih omejitev neposrednega praktičnega pomena.

Glavni cilj tega magistrskega dela je zgraditi in strogo ovrednotiti pristop k reševanju ONP, ki napredne transformerske modele za oceno momentov donosov združuje z metahevristično optimizacijo ob kardinalni omejitvi. Kot protiutež šibkim napovedim vpeljemo dodatno raven robustnosti --- dinamično združevanje transformerskih in zgodovinskih napovedi s postopkom Hedge~\cite{freund1997}, ki zagotavlja, da kombinacija v nobenem obdobju ne zaostane bistveno za najboljšim posameznim virom informacij.

Delo je zgrajeno takole. Začeli bomo s pregledom področja (Poglavje~\ref{ch:ozadje}), kjer bomo formalno opredelili osnovne pojme, mere uspešnosti in kardinalno omejeni mean-variance problem. Nato bomo v Poglavju~\ref{ch:model} predstavili napovedni model --- presečni transformer, ki iz iste predstavitve napove donose in kovarianco --- ter spletni ansambel napovedi. V Poglavju~\ref{ch:algoritmi} bomo opisali tri metahevristike, s katerimi rešujemo kardinalno omejeni problem. Sledilo bo eksperimentalno vrednotenje (Poglavje~\ref{ch:eksperiment}), v katerem bomo celoten cevovod postopno preizkusili v šestih tržnih režimih in na protokolih dveh objavljenih raziskav. Delo bomo sklenili z razpravo o rezultatih in njihovih omejitvah (Poglavje~\ref{ch:razprava}) ter sklepom (Poglavje~\ref{ch:sklep}).

Posebno pozornost bomo namenili poštenemu vrednotenju: donose bomo merili postopno (angl.\ \emph{walk-forward}), tako da model nikoli ne vidi prihodnjih podatkov, statistično značilnost pa preverili z vrsto robustnih testov, vključno z na izbiro popravljenim Sharpovim količnikom (angl.\ \emph{Deflated Sharpe Ratio}), naš pristop pa neposredno primerjali tudi na protokolih nedavno objavljenih študij (Ramos in sod.~\cite{ramos2023}; Fernandes in Desell~\cite{fernandes2026}). S tem želimo pokazati, da lahko sodoben napovedni model ob skrbnem vrednotenju in robustnem združevanju napovedi klasični zgodovinski pristop k optimizaciji portfelja oprijemljivo \emph{izboljša} --- ne nadomesti ---, pridobljeno prednost pa ohrani tudi v neugodnih tržnih režimih.
